% Self-contained fragment for the research note; paper proof, not Lean-checked.
\section{A necessary correction to the countable-reservoir bridge}
\paragraph{Evidence status.}
The next two propositions are PROVED-ON-PAPER in this run. Their stochastic
statements are not Lean-checked. Exact rational finite-prefix checks are reported
separately and do not certify the infinite-horizon calculation. The finite-coin
Tier 3 theorem is unchanged and is not contradicted.

\begin{proposition}[One-law finite work does not transfer termination]
Consider fresh iid coins with deterministic gap $D=1/2$. Under $P$ their positive
sign probability is $5/8$, and under $Q$ it is $3/8$. Query each coin once. If
$X_i$ is its observed bit, stop at
\[
 T=\inf\{n\ge1:\ \textstyle\sum_{i=1}^n(2X_i-1)=1\}.
\]
Then $P(T<\infty)=1$, $\mathbb E_P T=4$, whereas
$Q(T<\infty)=3/5$. In particular $Q(T=\infty)=2/5$.
\end{proposition}
\begin{proof}
For iid bits of parameter $p$, put $q=1-p$. A first hit at time $2k+1$ has
$k+1$ upward and $k$ downward steps, and the first $2k$ partial sums never
exceed zero. The number of such paths is the Catalan number
$C_k=(k+1)^{-1}\binom{2k}{k}$. This follows by reflecting at the first
forbidden positive step: among paths with $k$ up and $k$ down steps, the
forbidden paths are counted by $\binom{2k}{k+1}$, whose subtraction gives
$C_k$. Thus
\[
 \Pr_p(T=2k+1)=C_kp^{k+1}q^k.
\]
The generating series $C(z)=\sum_{k\ge0}C_kz^k$ satisfies
$C(z)=1+zC(z)^2$: split a nonempty Catalan path at its first return to zero.
The solution continuous at zero with $C(0)=1$ is
$C(z)=(1-\sqrt{1-4z})/(2z)$ for $0\le z<1/4$.
Consequently the probability generating function, excluding the mass at
infinity, is
\[
 G_p(s)=\mathbb E_p[s^T;T<\infty]
       =psC(pqs^2)=\frac{1-\sqrt{1-4pqs^2}}{2qs},\qquad 0<s\le1.
\]
For $p\ne1/2$ all series converge in a neighborhood of $s=1$.
It follows that $G_p(1)=1$ when $p>1/2$, and $G_p(1)=p/q$ when $p<1/2$.
For $p>1/2$, differentiation gives $G_p'(1)=1/(2p-1)$.
Equivalently differentiate $G_p(s)=s[p+qG_p(s)^2]$ at $s=1$;
nonnegative-series differentiation computes $\mathbb E_pT$ because
$G_p(1)=1$. Substituting $p=5/8$ and $p=3/8$ gives the claims.
\end{proof}

The example lies exactly in the symmetric prior family $1/2\pm\eta$ with
$\eta=1/8$. It also has an interior-gap variant: choose $D=1/4$, retain the
same sign priors, and use only the first observation of each fresh coin.
The observed bits then have parameters $9/16$ and $7/16$; the same calculation
gives $\mathbb E_PT=8$ and $Q(T=\infty)=2/9$. Thus deterministic coin
endpoints are not responsible for the failure. For a fixed finite coin family,
the two latent-prior laws have a globally bounded likelihood ratio; for an
unbounded reservoir that finite-family comparison is unavailable.

This example still has finite full-transcript relative entropy. On every finite
stopped path the upward step count exceeds the downward count by one, so the
likelihood ratio $dP/dQ$ is $5/3$. On infinite paths it is zero, since P has no
such mass. Therefore $\operatorname{KL}(P,Q)=\log(5/3)$. It is consistent with
the proposed full-transcript bound at $h=1/2$: convexity of $1/x$ gives
$\log(5/3)\le8/15<128/225=(128/225)h^2\mathbb E_PT$.
Thus full KL can extend while the two-law stopping transfer cannot.

\begin{proposition}[Finite-prefix information implies terminal testing]
Let $P,Q$ describe a common policy with independent seed included in increasing
observable prefix sigma fields $\mathcal F_n$. Let $T$ be its stopping time.
Assume $T<\infty$ almost surely under each law and $\mathbb E_PT\le B<\infty$.
Suppose, for a constant $k\ge0$, every finite-prefix transcript satisfies
\[
 \operatorname{KL}(P|_{\mathcal F_n},Q|_{\mathcal F_n})
       \le k\mathbb E_P(T\wedge n).
\]
If a terminal test $\phi\in\{0,1\}$ has
$P(\phi=1)\ge1-\delta$ and $Q(\phi=1)\le\delta$, with $0<\delta<1/2$, then
\[
 \operatorname{kl}(1-\delta,\delta)\le kB.
\]
\end{proposition}
\begin{proof}
The event $A_n=\{T\le n,\phi=1\}$ belongs to $\mathcal F_n$ because the
test is measurable at stopping. Let $p_n=P(A_n)$ and $q_n=Q(A_n)$.
Then $p_n\ge1-\delta-P(T>n)\ge1-\delta-B/n$, while $q_n\le\delta$.
For all sufficiently large $n$, $1-\delta-B/n>\delta$.
Data processing for the binary indicator of $A_n$, and monotonicity of binary
relative entropy in its first argument above its second, and in its second
argument below its first, give
\[
 \operatorname{kl}(1-\delta-B/n,\delta)
 \le\operatorname{kl}(p_n,q_n)
 \le\operatorname{KL}(P|_{\mathcal F_n},Q|_{\mathcal F_n})\le kB.
\]
Continuity at $(1-\delta,\delta)$ proves the result by $n\to\infty$.
The case $q_n=0<p_n$ gives infinite binary relative entropy and cannot occur
under a finite information bound. No transfer of termination from $P$ to $Q$
has been used. Separate finite expected-work bounds under the two laws are
a sufficient condition for the stated two termination assumptions.
\end{proof}

\paragraph{Canonical numbering and the remaining formal obligation.}
After $n$ charged observations, at most $n$ coin identities have been observed.
Assign indices $1,\ldots,n$ in order of first observation and retain a dictionary
for revisits. Every unobserved iid latent draw can be deferred until its first
observation: prior decisions involving its label contain no information about
that draw. This recursive construction preserves the observed finite-prefix law
and the number of charged observations. It requires that each action completes
after finite computation; infinitely many cost-free operations do not constitute
an observation oracle. For the symmetric information bound, condition on the
first $n$ iid magnitudes and the independent seed, apply the deterministic-gap
finite-prefix bound, and integrate. Magnitudes may be appended for this upper
bound on information and then removed by data processing; the policy itself
does not observe them. For the rare-prior bound, the appropriate single-coin
inequality and bounded-ratio hybrid decomposition must instead be used.

This is a paper-level reservoir realization and reduction, not an existing Lean
theorem. The existing finite-family declaration cannot be instantiated with a
different coin count at each horizon and thereby produce a single already
formalized infinite-reservoir process. One must formalize its joint law and
prove consistent realization of every prefix. The terminal-testing proposition
then closes the limit once those prefix bounds and the policy's admissibility
are available. The minimax experiment's uniform expected budget gives finite
work under both alternatives, so the counterexample does not invalidate the
paper minimax lower bound.

\paragraph{Reusable scope and boundary.}
The terminal-testing bridge works for any observation channel and any positive
cost accounting once its finite-prefix KL bound has been proved. With unequal
costs, a proved bound of the form
$\operatorname{KL}_n\le c\mathbb E_P\sum_{j\le T\wedge n}a_j$
gives the corresponding testing bound by the same limiting argument, provided
stopping is established separately (a positive lower bound on $a_j$ is one
sufficient condition). There is no asserted universal Bernoulli scaling for
other channels. Cross-state information can invalidate the iid coin likelihood
factorization even though this abstract testing step still holds; a shared
unknown parameter learned at one state must be included in a different
finite-prefix information calculation.
